\section{Line Corner Detection}\label{line-corner-detection}

\subsection{Convolution Operator}\label{convolution-operator}

\subsubsection{Convolution VS
Cross-Correlation}\label{convolution-vs-cross-correlation}

\textbf{卷积}与\textbf{互相关}的核心思想相同：滑动一个核（kernel）在输入图像上，在每个位置计算对应元素乘积的和。

\paragraph{1D}\label{d}

卷积： \[
(f * g)[n] = \sum_{m=-\infty}^{\infty} f[m] \, g[n - m]
\] 互相关： \[
(f \star g)[n] = \sum_{m=-\infty}^{\infty} f[m] \, g[n + m]
\]

\paragraph{2D}\label{d-1}

卷积： \[
(f * g)[i, j] = \sum_{m=-\infty}^{\infty} \sum_{n=-\infty}^{\infty} f[i - m, j - n] \, g[m, n]
\] 互相关： \[
(f \star g)[i, j] = \sum_{m=-\infty}^{\infty} \sum_{n=-\infty}^{\infty} f[i + m, j + n] \, g[m, n]
\]

\paragraph{Difference}\label{difference}

卷积在滑动前会将核翻转（flip），在一维中是时间反转；互相关不翻转核。

\noindent\small
\setlength{\tabcolsep}{2pt}
\begin{tabular}{@{}lp{0.37\columnwidth}p{0.37\columnwidth}@{}}
\hline
 & \textbf{卷积 (Convolution)} & \textbf{互相关 (Correlation)} \\
\hline
\textbf{直观理解} & 滤波 (Filter) & 模板匹配 (Template matching) \\
\textbf{交换律} & Y & N \\
\textbf{结合律} & Y & N \\
\textbf{分配律} & Y & N \\
\hline
\end{tabular}

\subsubsection{Padding}\label{padding}

  防止深网络中 spatial
  shrinkage，保留边缘信息。

\begin{itemize}
\tightlist
\item
  Zero padding：以 0 填充。
\item
  Replicate padding：以边缘像素复制填充。
\end{itemize}

\textbf{Output Size}（各方向分别计算）：
\[
W_{out} = \left\lfloor \frac{W_{in} + 2P - K}{S} \right\rfloor + 1
\]
{\scriptsize $W_{in}$: 输入宽; $P$: 单侧填充; $K$: 核宽; $S$: 步幅; $W_{out}$: 输出宽}

\subsubsection{Implementation of
Convolution}\label{implementation-of-convolution}

用循环实现卷积在 CPU/GPU
上效率低。想办法将卷积转化为矩阵乘法并使用并行可以大大提高速率。

\paragraph{im2col}\label{im2col}

\textbf{符号定义}：

符号 \textbar{} 含义 \textbar{} 维度 \textbar{}\\
:--- \textbar{} :--- \textbar{} :--- \textbar{}\\
\(\mathbf{X} \in \mathbb{R}^{C_{in} \times H \times W}\) \textbar{}
输入特征图 \textbar{} \((C_{in}, H, W)\) \textbar{}\\
\(\mathbf{K} \in \mathbb{R}^{C_{out} \times C_{in} \times K \times K}\)
\textbar{} 卷积核权重 \textbar{} \((C_{out}, C_{in}, K, K)\)
\textbar{}\\
\(\mathbf{b} \in \mathbb{R}^{C_{out}}\) \textbar{} 偏置向量 \textbar{}
\((C_{out},)\) \textbar{}\\
\(\mathbf{Y} \in \mathbb{R}^{C_{out} \times H_{out} \times W_{out}}\)
\textbar{} 输出特征图 \textbar{} \((C_{out}, H_{out}, W_{out})\)
\textbar{}\\
\((s_h, s_w)\) \textbar{} 步长 \textbar{} - \textbar{}\\
\((p_h, p_w)\) \textbar{} 填充 \textbar{} - \textbar{}

其中输出空间尺寸： \[
H_{out} = \left\lfloor \frac{H + 2p_h - K}{s_h} \right\rfloor + 1,
\quad W_{out} = \left\lfloor \frac{W + 2p_w - K}{s_w} \right\rfloor + 1
\]

令 \(N = H_{out} \cdot W_{out}\) 表示输出位置的总数。

\textbf{Kernel Flattening}：

将每个输出通道的卷积核\textbf{展平为行向量}：

\[
\mathbf{W}_{i} = \text{vec}\left(\mathbf{K}_{i, :, :, :}\right)^\top \in \mathbb{R}^{1 \times (C_{in} \cdot K^2)}
\]

其中 \(\text{vec}(\cdot)\) 表示按通道优先顺序展平： \[
\text{vec}\left(\mathbf{K}_{i, :, :, :}\right)
= \left[ \text{vec}\left(\mathbf{K}_{i, 0, :, :}\right)^\top, \text{vec}\left(\mathbf{K}_{i, 1, :, :}\right)^\top,
\ldots, \text{vec}\left(\mathbf{K}_{i, C_{in}-1, :, :}\right)^\top \right]^\top
\]

堆叠所有输出通道形成\textbf{权重矩阵}： \[
\mathbf{W} = \begin{bmatrix} \mathbf{W}_0 \\ \mathbf{W}_1 \\ \vdots \\ \mathbf{W}_{C_{out}-1} \end{bmatrix}
\in \mathbb{R}^{C_{out} \times (C_{in} \cdot K^2)}
\]

\textbf{Patch Flattening}：

对于每个输出位置 \((i, j)\)，\(i \in [0, H_{out})\),
\(j \in [0, W_{out})\)，定义对应的输入感受野：

\[
\mathcal{P}_{i,j} = \left\{ (c, h, w) \,\middle|\, c \in [0, C_{in}),\, h \in [h_0, h_0+K),\, w \in [w_0, w_0+K) \right\}
\] 其中 \((h_0, w_0) = (i \cdot s_h - p_h, j \cdot s_w - p_w)\)。

提取 patch 并\textbf{展平为列向量}： \[
\mathbf{x}_{i,j} = \text{vec}\left(\mathbf{X}_{:, h_0:h_0+K, w_0:w_0+K}\right) \in \mathbb{R}^{C_{in} \cdot K^2}
\]

\textbf{im2col}：

将所有输出位置对应的列向量按列拼接，形成 \textbf{im2col 矩阵}： \[
\mathbf{X}_{col} =
\begin{bmatrix} \mathbf{x}_{0,0} & \cdots & \mathbf{x}_{0,W_{out}-1} & \mathbf{x}_{1,0} & \cdots
& \mathbf{x}_{H_{out}-1,W_{out}-1} \end{bmatrix} \in \mathbb{R}^{(C_{in} \cdot K^2) \times N}
\]

其中列索引 \(t = i \cdot W_{out} + j\) 对应输出位置 \((i, j)\)。

\textbf{GEMM}：

执行矩阵乘法并加上偏置： \[
\mathbf{Y}_{mat}
= \mathbf{W} \cdot \mathbf{X}_{col} + \mathbf{b} \cdot \mathbf{1}_N^\top \in \mathbb{R}^{C_{out} \times N}
\]

其中 \(\mathbf{1}_N \in \mathbb{R}^{N}\) 为全 1
向量，偏置通过广播机制相加。

\textbf{Reconstruction}：

将结果矩阵 reshape 回张量形式：

\[
\mathbf{Y} = \text{reshape}\left(\mathbf{Y}_{mat}, (C_{out}, H_{out}, W_{out})\right)
\]

具体地，对于每个输出通道 \(c_{out}\) 和空间位置 \((i, j)\)： \[
\begin{aligned}
\mathbf{Y}_{c_{out}, i, j}
&= \left[\mathbf{Y}_{mat}\right]_{c_{out},\, i \cdot W_{out} + j} \\
&= \sum_{d=0}^{C_{in} \cdot K^2 - 1} \mathbf{W}_{c_{out}, d} \cdot
\left[\mathbf{X}_{col}\right]_{d, i \cdot W_{out} + j} + \mathbf{b}_{c_{out}} \\
&= \mathbf{W}_{c_{out}} \cdot \mathbf{x}_{i, j} + \mathbf{b}_{c_{out}}
\end{aligned}
\]

\textbf{全过程}：

\[
\mathbf{X}
\xrightarrow{\text{im2col}} \mathbf{X}_{col}
\xrightarrow{\mathbf{W} \cdot (\cdot) + \mathbf{b}} \mathbf{Y}_{mat}
\xrightarrow{\text{reshape}} \mathbf{Y}
\]

\subsection{Line Fitting}\label{line-fitting}

\subsubsection{Least Squares}\label{least-squares}

若已知一系列点 \((x_i, y_i)\)，希望拟合直线
\(y = kx + m\)。可以用最小二乘法通过最小化残差平方和求解参数。 也即： \[
\min_{k,m} \sum_{i=1}^n (y_i - (k x_i + m))^2
\] 矩阵形式： \[
\min_{\boldsymbol{\theta}} \| A\boldsymbol{\theta} - \boldsymbol{y} \|_2^2
\Leftrightarrow \min_{\boldsymbol{\theta}} (A\boldsymbol{\theta} - \boldsymbol{y})^T (A\boldsymbol{\theta} - \boldsymbol{y})
\] 其中 \[
A = \begin{bmatrix} x_1 & 1 \\ x_2 & 1 \\ \vdots & \vdots \\ x_n & 1 \end{bmatrix},\quad
\boldsymbol{\theta} = \begin{bmatrix} k \\ m \end{bmatrix},\quad
\boldsymbol{y} = \begin{bmatrix} y_1 \\ y_2 \\ \vdots \\ y_n \end{bmatrix}
\] 最小二乘解析解： \[
\boldsymbol{\theta} = (A^\top A)^{-1} A^\top \boldsymbol{y}
\] \textbf{注}： 若拟合直线接近竖直，原有斜率参数不稳定。

\subsubsection{Imperoved Least Squares}\label{imperoved-least-squares}

由于前述方法下，当直线竖直时，斜率无穷，最小二乘解不稳定。故我们应当使用一般直线方程求解。也即：
\[
\min_{a,b,d} \sum_{i=1}^n (a x_i + b y_i - d)^2
\] 但是为了解的唯一性（并且防止平凡解
\(a = b = d = 0\)），我们可以添加约束
\(a^2 + b^2 + d^2 = 1\)。得到矩阵形式： \[
\min_{\boldsymbol{h}} \|A \boldsymbol{h}\|^2
\] 其中 \[
A = \begin{bmatrix} x_1 & y_1 & 1 \\ x_2 & y_2 & 1 \\ \vdots & \vdots & \vdots \\ x_n & y_n & 1 \end{bmatrix},\quad
\boldsymbol{h} = \begin{bmatrix} a \\ b \\ d \end{bmatrix},\quad
\|\boldsymbol{h}\| = 1
\]

我们可以通过 \textbf{SVD} 求解。对矩阵 \(A\) 进行奇异值分解得： \[
A_{n\times 3} = U_{n\times n} D_{n\times 3} V_{3\times 3}^\top
\] 其中 \(U\) 为 \(n \times n\)
正交阵，\(V = \begin{bmatrix} \boldsymbol{v}_1^T \\ \boldsymbol{v}_2^T \\ \boldsymbol{v}_3^T \end{bmatrix}\)
为 \(3 \times 3\)
正交阵，\(D = \begin{bmatrix} diag\{\lambda_1, \lambda_2, \lambda_3\} \\ O_{(n-3) \times 3} \end{bmatrix}\)
为类对角矩阵，包含奇异值
\(|\lambda_1| \ge |\lambda_2| \ge |\lambda_3|\)。 将 \(h\) 分解到 \(V\)
的正交标架下： \[
\boldsymbol{h} = \alpha_1 \boldsymbol{v}_1 + \alpha_2 \boldsymbol{v}_2 + \alpha_3 \boldsymbol{v}_3
\] 由于 $ \textbar{}\boldsymbol{h}\textbar{} = 1 $ 从而
\(\alpha_1^2 + \alpha_2^2 + \alpha_3^2 = 1\) ： \[
\begin{aligned}
\| Ah \| = \| U D V^T h \| = \| D V^T h \| \\
= \| D \begin{bmatrix} \alpha_1 \\ \alpha_2 \\ \alpha_3 \end{bmatrix} \|
= \| \begin{bmatrix} \lambda_1 \alpha_1 \\ \lambda_2 \alpha_2 \\ \lambda_3 \alpha_3 \\ O^T \end{bmatrix} \|
= (\lambda_1 \alpha_1)^2 + (\lambda_2 \alpha_2)^2 + (\lambda_3 \alpha_3)^2
\end{aligned}
\] 从而解析解为 \(\boldsymbol{h} = \pm \boldsymbol{v}_3\).
几何直观即：解为最小伸缩的方向。

\subsubsection{RANSAC}\label{ransac}

最小二乘对异常值（outliers）非常敏感，少数严重偏离点会破坏整体拟合。

随机采样一致性算法（\textbf{RAN}dom \textbf{SA}mple
\textbf{C}onsensus）可以解决这个问题。
思想是通过大量随机采样最小样本集，生成候选模型并统计内点，最终选择支持内点最多的模型。

\paragraph{Step of RANSAC}\label{step-of-ransac}

\begin{enumerate}
\def\labelenumi{\arabic{enumi}.}
\tightlist
\item
  确定构成模型所需的最小样本数 \(s\)、残差阈值 \(\delta\)、期望成功概率
  \(p\)（或最大迭代次数 \(N\)）。
\item
  从数据集中随机选择 \(s\) 个样本。
\item
  根据所选样本拟合一个假设模型（putative model）。
\item
  计算所有点相对于该模型的残差（residual），并判断哪些点为内点（\(residual < \delta\)）。
\item
  若当前模型的内点数历史最佳，则更新模型与内点集。
\item
  重复上述步骤 \(N\)
  次。结束后，用内点集合重新拟合得到最终模型（可用最小二乘或 SVD）。
\end{enumerate}

\paragraph{The Number of Sample}\label{the-number-of-sample}

上述步骤中，选择合适的迭代次数 \(N\) 是很重要的。

假设数据集整体服从一个真实模型。我们应当使得完成迭代后，出现一次优秀抽样（抽的
\(s\)
个样本全是真实模型的内点）的概率尽量大。假设构成模型所需的最少样本数为
\(s\)，真实模型中外点比例为 \(e\)，期望出现优秀抽样的概率为 \(p\)。

则一次抽样为非优秀抽样（ \(s\) 个样本中存在真实模型外点）的概率： \[
1 - (1 - e)^s
\] \(N\) 次抽样中不存在优秀抽样的概率： \[
(1 - (1 - e)^s)^N = 1 - p
\] 解得： \[
N = \frac{\log(1 - p)}{\log(1 - (1 - e)^s)}
\]

\paragraph{Imperovement}\label{imperovement}

\begin{enumerate}
\def\labelenumi{\arabic{enumi}.}
\tightlist
\item
  对于超参数 threshold 的选取很重要，可以动态调整进行选取。
\item
  由于对噪声的敏感性，仅做一次 RANSAC
  不一定能建立很好的模型，但是可以排除很多 outliers。我们可以对剩下的
  inliers 再做 Line Fitting（如使用 Least Squares 或者 SVD）.
\end{enumerate}

\subsection{Corner Detection}\label{corner-detection}

\subsubsection{Harris Detector}\label{harris-detector}

思想：若窗口向任意方向平移，窗口内像素强度都会显著变化，则该点可能为角点。

\paragraph{The Definition fo Energy}\label{the-definition-fo-energy}

考虑一个以 \((x_0,y_0)\) 为中心的窗口 \(N(x_0,y_0)\)，平移
\((u,v)\)，窗口内像素强度的变化程度定义为该点处的\textbf{能量}： \[
E_{x_0,y_0}(u,v) = \sum_{(x,y)\in N(x_0,y_0)} [I(x+u,y+v) - I(x,y)]^2
\] 其中 \(I(\cdot,\cdot)\) 为图像强度。 使用一阶泰勒展开（假设
\(u, \; v\) 很小）可得： \[
I(x+u,y+v) - I(x,y) \approx I_x(x,y) u + I_y(x,y) v
\] 代回能量： \[
E(u,v) \approx \sum_{(x,y)\in N} (I_x(x,y) u + I_y(x,y) v)^2
\] 利用二次型化简得： \[
E(u,v) \approx \begin{bmatrix} u & v \end{bmatrix}
M(x_0,y_0)
\begin{bmatrix} u \\ v \end{bmatrix}
\] 其中 \[
M(x_0,y_0)
= \sum_{(x,y)\in N(x_0,y_0)} \begin{bmatrix} I_x^2(x,y) & I_x(x,y) I_y(x,y) \\ I_x(x,y) I_y(x,y) & I_y^2(x,y) \end{bmatrix}
\]

\paragraph{Window Function}\label{window-function}

我们可以进一步将二次型 \(M\) 表示为卷积形式。常见的有以下两种窗口：

\begin{itemize}
\item
  Rectangle window：窗函数对窗内均匀加权，不具旋转不变性。

  \[
  M(x_0,y_0)
  =
  \begin{bmatrix}
  w(x_0,y_0) * I_x^2 & w(x_0,y_0) * I_x I_y \\
  w(x_0,y_0) * I_x I_y & w(x_0,y_0) * I_y^2
  \end{bmatrix}
  \] 其中
  \(w(x_0, y_0) = \begin{cases} 1, \quad if (x, y) \in N(x_0, y_0) \\ 0, \quad otherwise\end{cases}\)
\item
  Gaussian window：用高斯核加权，具有旋转不变性和平滑性。

  \[
  M(x_0,y_0)
  =
  \begin{bmatrix}
  g_\sigma(x_0,y_0) * I_x^2 & g_\sigma(x_0,y_0) * I_x I_y \\
  g_\sigma(x_0,y_0) * I_x I_y & g_\sigma(x_0,y_0) * I_y^2
  \end{bmatrix}
  \] 其中高斯窗口： \[
  g_\sigma(x_0,y_0) = \frac{1}{2\pi\sigma^2} \exp \big(-\frac{x_0^2+y_0^2}{2\sigma^2}\big)
  \]
\end{itemize}

当然由于卷积是一种线性变换，我们可以统一写为： \[
M(x_0,y_0)
=
\begin{bmatrix}
\mathcal{F}^{(x_0, y_0)}(I_x^2) & \mathcal{F}^{(x_0, y_0)}(I_x I_y) \\
\mathcal{F}^{(x_0, y_0)}(I_x I_y) & \mathcal{F}^{(x_0, y_0)}(I_y^2)
\end{bmatrix}
\] 其中
\(\mathcal{F}^{(x_0, y_0)}: \mathbb{R}^{n \times n} \rightarrow \mathbb{R}\)

\paragraph{Eigenvalues for
Classification}\label{eigenvalues-for-classification}

\(M\) 是对称正半定矩阵，可做特征值分解： \[
M = Q \begin{bmatrix} \lambda_1 & 0 \\ 0 & \lambda_2 \end{bmatrix} Q^\top
\] 其中 \(\lambda_1, \lambda_2 \ge 0\) 为 \(M\) 的特征值。 由此得： \[
E_{(x_0, y_0)}(u, v) \approx \lambda_1 u'^2 + \lambda_2 v'^2
\] 其中
\(\begin{bmatrix} u' \\ v' \end{bmatrix} = Q \begin{bmatrix} u \\ v \end{bmatrix}\).

由此可见 \(E(u, v)\)
是一个抛物面，两个特征值分别对应抛物面在主轴方向上的曲率。依据
\(\lambda_1, \lambda_2\) 的大小可判定点的类型：

\begin{itemize}
\tightlist
\item
  平坦区域（Flat）：能量在任方向都很小，也即
  \(\lambda_1 \approx 0, \lambda_2 \approx 0\)。
\item
  边缘（Edge）：沿边缘方向平移能量变化小，垂直方向平移变化大，也即一个特征值大，另一个接近
  0。
\item
  角点（Corner）：沿任意方向平移都会引起能量较大变化，也即两个特征值都大。
\end{itemize}

\paragraph{Corner Response Function}\label{corner-response-function}

上述判断需要做特征分解，计算成本略高。我们使用一个近似响应函数用以判断角点。
希望角点满足： \[
1/k < \lambda_1/\lambda_2 < k \qquad \tag{1}
\] \[
\lambda_1, \lambda_2 > b \qquad \tag{2}
\] 令 \[
\begin{aligned}
\theta
&= \frac{1}{2} \underbrace{(\lambda_1 \lambda_2 - 2 \alpha (\lambda_1 + \lambda_2)^2)}_{(1)}
+ \frac{1}{2} \underbrace{(\lambda_1 \lambda_2 - 2t)}_{(2)} \\
&= \lambda_1 \lambda_2 - \alpha (\lambda_1 + \lambda_2)^2 - t \\
&= \det{M} - \alpha (\mathbb{tr}M)^2 - t \\
&= (\mathcal{F}(I_x^2) \mathcal{F}(I_y^2) - \mathcal{F}(I_x I_y)^2)
- \alpha (\mathcal{F}(I_x^2) + \mathcal{F}(I_y^2))^2 - t
\end{aligned}
\] 即为 Harris 响应。若 \(\theta\) 大且为正则为角点；若 \(\theta\)
负并且绝对值大则为边缘；若 \(\theta\) 接近 0 则为平坦区域。

\textbf{注}：其中 \(\alpha, \; t\) 均为超参数。

\paragraph{Step of Harris}\label{step-of-harris}

Harris Detector 在实际实现中通常步骤为：

\begin{enumerate}
\def\labelenumi{\arabic{enumi}.}
\tightlist
\item
  计算图像梯度 \(I_x, I_y\)，并计算
  \(I_x^2, I_y^2, I_x I_y\)，得到与原图像同维度的三张梯度图。
\item
  对上述三图用高斯滤波 \(g_\sigma\) 卷积得到三张与原图像同维度的高斯图。
\item
  对于像素点
  \((x_0, y_0)\)，使用三张高斯图对应位置的三个数构成二阶对称阵，计算
  \(\theta\)，并进行阈值二值化 \(\theta > 0\)。
\item
  非极大值抑制（\textbf{NMS}） 保留局部极大值作为关键点位置。
\end{enumerate}

\paragraph{Equivariance \& Invariance}\label{equivariance-invariance}

设输入信号空间 \(\mathbb{V}\)，输出表示空间 \(\mathbb{W}\)，信号检测函数
\(f: \mathbb{V} \rightarrow \mathbb{W}\)。

记 \(\mathcal{T}\) 为抽象变换集合（如平移 5 像素、旋转
\(\pi/2\)）。对每个 \(\tau \in \mathcal{T}\)，同时定义：

\begin{itemize}
\tightlist
\item
  \(T_{\mathbb{V}}^{(\tau)}: \mathbb{V} \rightarrow \mathbb{V}\)（\(\tau\)
  在输入空间的实现）
\item
  \(T_{\mathbb{W}}^{(\tau)}: \mathbb{W} \rightarrow \mathbb{W}\)（\(\tau\)
  在输出空间的实现）
\end{itemize}

\textbf{Equivariance}（等变性）：对输入信号做变换，输出表示也同样变换。
\[
T_{\mathbb{W}}^{(\tau)}[f(X)] = f(T_{\mathbb{V}}^{(\tau)}(X)), \quad \forall \tau \in \mathcal{T}
\]

\textbf{Invariance}（不变性）：对输入信号做变换，输出表示不变。 \[
f(X) = f(T_{\mathbb{V}}^{(\tau)}(X)), \quad \forall \tau \in \mathcal{T}
\]

\textbf{注}：上述性质都是检测函数 \(f\) 的。

\textbf{Harris Detector 的性质}：

记输入图像为 \(A \in \mathbb{R}^{n \times n}\)。

\[
M(x_0,y_0)
=
\begin{bmatrix}
\mathcal{F}^{(x_0, y_0)}(I_x^2) & \mathcal{F}^{(x_0, y_0)}(I_x I_y) \\
\mathcal{F}^{(x_0, y_0)}(I_x I_y) & \mathcal{F}^{(x_0, y_0)}(I_y^2)
\end{bmatrix}
\]

\[
\theta
= (\mathcal{F}(I_x^2) \mathcal{F}(I_y^2) - \mathcal{F}(I_x I_y)^2)
- \alpha (\mathcal{F}(I_x^2) + \mathcal{F}(I_y^2))^2 - t
\]

\[
\Theta(A) = [\theta(i, j)] \in \mathbb{R}^{n \times n}
\]

以下的讨论不考虑边缘处。

\begin{itemize}
\item
  \textbf{所有}的 \(\Theta\) 对图像\textbf{平移（Translation）等变}：

  记 \(T_{u, v}\) 表示平移 \((u, v)\) 个像素，也即
  \((T_{u,v}(A))_{x, y} = A_{x-u, y-v}\)。则有： \[
  T_{u, v}(\Theta(A)) = \Theta(T_{u, v}(A))
  \]
\item
  \textbf{各向同性（Isotropic）} 的 filter 产生的 \(\Theta\)
  对\textbf{图像旋转（Rotation）等变}：

  记 \(R_\phi\) 表示图像旋转 \(\phi\) 角度，即
  \((R_\phi(A))_{x,y} = A_{x',y'}\)，其中 \((x', y')\) 是 \((x,y)\) 旋转
  \(-\phi\) 后的坐标。 filter 为各向同性的核 \(g\)。 则有： \[
  R_\phi(\Theta(A)) = \Theta(R_\phi(A))
  \]

  \textbf{注}：

  \begin{itemize}
  \tightlist
  \item
    约等于是因为在像素网格中考虑，旋转后必须采用插值方式计算像素值。
  \item
    kernel 的 rotation \textbf{invariant} 对应于 kernel convolution 的
    rotation \textbf{equivariant}。
  \end{itemize}
\item
  对\textbf{尺度（Scale）不等变}： 放大缩小 viewpoint
  会导致角点检测结果变化。
\end{itemize}
